\section{Relative primary signatures and specialization laws}
\label{sec:relative-primary-specialization}

This section supplies the two family arguments that are not consequences
of a single exact witness: genericity on coefficient orbits and
fibrewise control of associated points.  It also extracts the
geometric primary signatures of the corank-two layers and writes
explicit orbit degenerations.

\subsection{Stabilizer slices and generic Groebner bases}

Keep
\[
 P=\C[a,b,c,d],\qquad
 T=\begin{pmatrix}a&b\\c&d\end{pmatrix},\qquad
 \delta=ad-bc.
\]
For a mixed normal form \(\beta\), let
\(G_\beta\subset GL(H)\times GL(W)\times GL(E_H)\) be its
stabilizer.  If \((X,Y,Z)\) lies in its Lie algebra, then
\begin{equation}\label{eq:stabilizer-linearization}
 Z\beta-\beta(X\otimes1+1\otimes Y)=0,
\end{equation}
and its infinitesimal action on the quadratic block is
\begin{equation}\label{eq:chi-linearization}
 \dot\chi=Z\chi-\chi\,\Sym^2Y.
\end{equation}
These are linear equations over \(\Q\) at the integral base
points of Lemma~\ref{lem:nine-orbit-colons}.

\begin{theorem}[Generic slice certification for the mixed-kernel atlas]
\label{thm:generic-slice-certification}
For each of the seven mixed-kernel rows not already covered by the
mixed-regular and decomposable-kernel theorems, the Hilbert functions
in Theorem~\ref{thm:corank-two-mixed-kernel-atlas} are the generic
Hilbert functions on the entire orbit family.

More precisely, at the integral base points used in
Lemma~\ref{lem:nine-orbit-colons}, the ranks of the infinitesimal
\(G_\beta\)-orbits in the nine-dimensional \(\chi\)-space and
one choice of transverse affine slice are
\[
\begin{array}{c|c|c}
\text{type}&\rank T_{\chi}(G_\beta\chi)&\text{slice entries}\\ \hline
\text{secant}&6&\chi_{11},\chi_{21},\chi_{31}\\
\text{tangent}&7&\chi_{11},\chi_{21}\\
H\text{-ruling}&7&\chi_{11},\chi_{13}\\
W\text{-ruling}&9&\varnothing\\
\text{dual rank }2&9&\varnothing\\
\text{dual rank }1&9&\varnothing\\
\text{zero mixed map}&9&\varnothing.
\end{array}
\]
On the first three slices, with coordinates \(s_i\), Groebner
bases of all required colons over \(\Q(s_i)\) give respectively
\[
\begin{array}{c|c|c|c}
\text{type}&\min\{k:\delta^k\in J\}&H_{W_2}&H_{W_3}\\ \hline
\text{secant}&2&(1,4,3)&0\\
\text{tangent}&2&(1,4,3)&0\\
H\text{-ruling}&3&(1,4,6,4)&(1).
\end{array}
\]
The only nonconstant denominators in reduced bases may be chosen to
divide
\[
 s_2-2\quad\text{(secant)},\qquad
 s_0+2s_1-5\quad\text{(\(H\)-ruling)},
\]
and neither vanishes at the base point \(s=0\); the tangent slice
requires no nonconstant denominator.  In the remaining four rows the
stabilizer orbit is already open in \(\chi\)-space, so the exact
base-point calculation is the generic calculation.
\end{theorem}

\begin{proof}
Solving \eqref{eq:stabilizer-linearization} and applying
\eqref{eq:chi-linearization} gives the ranks in the table.  Adjoining
the displayed coordinate directions raises the tangent rank to nine.
Thus the morphism from the stabilizer times the affine slice to
\(\chi\)-space has surjective differential at the base point.
Over characteristic zero it is therefore dominant and smooth on a
neighbourhood of that point.

Compute the residual minors and successive colons in
\(\Q(s_i)[a,b,c,d]\).  For the secant and tangent slices the
second-layer generic initial ideal is
\[
 (c,d)^3+(a,b)^2+(a,b)(c,d),
\]
which has Hilbert function \((1,4,3)\), and the next colon is the
unit ideal.  For the \(H\)-ruling slice a generic initial ideal of
the second layer is generated by
\[
 c^4,c^3d,c^2d^2,cd^3,d^4,
 ac^2,acd,ad^2,bd^2,a^2,ab,b^2,bc,
\]
with Hilbert function \((1,4,6,4)\); the third-layer initial ideal
is \((a,b,c,d)\).  These initial ideals are unchanged after
specialization away from the displayed denominator divisors.

For the four open-orbit rows the exact Groebner bases give the
functions printed in Theorem~\ref{thm:corank-two-mixed-kernel-atlas}.
The calculation is reproduced, over exact rational function fields
rather than floating-point samples, by
\texttt{checks/generic\_boundary\_atlas.py}.  Since the slice map
is dominant, these are the generic initial ideals on the orbit
families.
\end{proof}

\subsection{Geometric primary signatures for the nine mixed-kernel types}

Let \(Q=V(\delta)\) be the affine Segre cone and let
\(\mathfrak m=(a,b,c,d)\) be its vertex.  Write \(R_r\) for a
reduced divisor consisting of \(r\) distinct lines in one ruling of
\(Q\).  Write \(V_\ell\) for a homogeneous
\(\mathfrak m\)-primary scheme of length \(\ell\).  Finally,
\(R_r+E_\ell\) means a scheme whose saturation is \(R_r\) and
whose saturation quotient is an embedded vertex module of length
\(\ell\).  This notation records associated supports and generic
lengths; it does not choose a noncanonical embedded primary
component.

\begin{theorem}[Generic primary signatures for the mixed-kernel orbits]
\label{thm:generic-primary-signatures}
On the generic coefficient stratum of each \(A_H\)-injective
mixed-kernel orbit, the positive layers have the following primary
signatures:
\[
\begin{array}{c|c|c}
\text{mixed-kernel type}&W_2&W_3\\ \hline
3,\ \text{point off }Q&V_1&\varnothing\\
3,\ \text{point on }Q&V_3&\varnothing\\
2,\ \text{secant}&V_8&\varnothing\\
2,\ \text{tangent}&V_8&\varnothing\\
2,\ H\text{-ruling}&V_{15}&V_1\\
2,\ W\text{-ruling}&R_1+E_5&\varnothing\\
1,\ \text{dual rank }2&R_2+E_3&V_1\\
1,\ \text{dual rank }1&R_3&V_3\\
0&Q&Q.
\end{array}
\]
In particular, every associated support and every generic length in
these layers is determined.  The symbols \(R_r\) have \(r\)
minimal ruling-line primes, each of generic length one.  The only
additional associated prime in \(R_1+E_5\) and
\(R_2+E_3\) is the embedded vertex \(\mathfrak m\).
\end{theorem}

\begin{proof}
The finite rows are homogeneous zero-dimensional schemes supported at
the vertex, hence are \(\mathfrak m\)-primary; their lengths are
the Hilbert-function sums.  For the \(W\)-ruling witness the second
layer is
\[
 I=(ac,bc,ad,bd,c^3,c^2d,cd^2,d^3).
\]
Its saturation is \((c,d)\), and \((c,d)/I\) has length five.
Thus its associated primes are \((c,d)\) and \(\mathfrak m\),
with an embedded vertex contribution of length five.

For the dual-rank-two witness, saturation gives
\[
 (a,b)\cap(c,d),
\]
the union of two ruling lines, and the saturation quotient has length
three.  For the dual-rank-one witness, modulo \(\delta\), the
second-layer ideal is exactly
\[
 \bigcap_{\lambda\in\{0,1,3\}}
 (a+\lambda c,b+\lambda d),
\]
so it is a reduced union of three ruling lines and has no embedded
point.  The third layer is the length-three vertex thickening already
computed in Lemma~\ref{lem:nine-orbit-colons}.  The zero mixed-map
identity \(J=(\delta^3)\) gives \(W_2=W_3=Q\).

Theorem~\ref{thm:generic-slice-certification} transports these exact
signatures to the generic coefficient strata.  Associated supports
are preserved by the stabilizer action, and the relative-primary
argument below prevents an isolated base-point calculation from being
used beyond its certified open stratum.
\end{proof}

\subsection{Pure-block rank drop at projection corank two}

When \(A_H\) has rank \(a<3\), put \(q=6-a\) and
\(b=\rank\overline B_H\).  For \(b=3,2,1\) the phrase
\emph{generic mixed kernel} means respectively an off-Segre kernel
point, a secant kernel line, and a kernel plane whose dual tensor has
rank two.

\begin{theorem}[Generic primary signatures for every \(A_H\)-rank-drop type]
\label{thm:rankdrop-primary-signatures}
At projection corank two, every generic rank type allowed by
Proposition~\ref{prop:rankdrop-coranktwo-depth} has the following
complete list of nonzero positive layers:
\[
\begin{array}{c c|c|l}
a&b&\text{nilpotency index}&(W_1,W_2,W_3,W_4,W_5)\\ \hline
2&4&3&(Q,V_{22})\\
2&3&4&(Q,V_{27},V_1)\\
2&2&4&(Q,R_4,V_8)\\
2&1&5&(Q,Q,Q,V_1)\\
1&4&5&(Q,Q,V_{11},V_1)\\
1&3&5&(Q,Q,R_2+E_3,V_1)\\
1&2&5&(Q,Q,Q,R_2)\\
0&4&6&(Q,Q,Q,R_2,V_1)\\
0&3&6&(Q,Q,Q,Q,V_1).
\end{array}
\]
Here every \(Q\) is the reduced Segre quadric with generic length
one; every \(R_r\) is a reduced union of \(r\) ruling lines with
generic length one on each component; and every \(V_\ell\) is
\(\mathfrak m\)-primary of the indicated length.  In the unique
row \((a,b)=(1,3)\) carrying \(E_3\), the associated primes of
that layer are the two ruling-line primes together with the embedded
vertex.
\end{theorem}

\begin{proof}
Choose the generic normal forms of the mixed block described above and
let \(\chi\) vary.  The stabilizer-slice codimensions in
\(\chi\)-space for
\[
 (a,b)=(2,4),(2,3),(2,2),(2,1),(1,4),(1,3),(1,2),(0,4),(0,3)
\]
are respectively
\[
 5,4,1,0,3,1,0,0,0.
\]
On explicit complementary coordinate slices the successive colons are
computed over the rational function fields of those slices.  Their
Hilbert functions are
\[
\begin{array}{c c|l}
a&b&\text{non-quadric Hilbert functions}\\ \hline
2&4&(1,4,9,8)\\
2&3&(1,4,9,8,5);\ (1)\\
2&2&4n+4\ (n\ge4);\ (1,4,3)\\
2&1&(1)\\
1&4&(1,4,6);\ (1)\\
1&3&2n+2\ (n\ge3);\ (1)\\
1&2&2n+2\ (n\ge1)\\
0&4&2n+2\ (n\ge1);\ (1)\\
0&3&(1).
\end{array}
\]
Every omitted earlier layer has Hilbert function \((n+1)^2\),
hence equals \(Q\).

To identify the positive-dimensional supports, pull the layer ideals
back to the rank-one parametrization \(T=uv^t\), but now perform the
factorization over the same rational-function fields as the Groebner
calculation.

For \((a,b)=(2,2)\), with slice parameter \(s\), the second layer has
common rank-one factor, up to a unit of \(\Q(s)\),
\[
 u_0u_1\Bigl(
 (s+9)u_0^2+(3s+6)u_0u_1-7u_1^2
 \Bigr).
\]
The quadratic discriminant is
\[
 9s^2+64s+288.
\]
On the open set
\((s+9)(9s^2+64s+288)\ne0\), its two roots are distinct and
different from the two roots \(u_0u_1=0\).  Thus the generic support
is a reduced divisor of four distinct ruling lines.  Its Hilbert
polynomial agrees with the computed layer, so there is no generic
embedded correction.

For \((a,b)=(1,3)\), with slice parameter \(t\), the third layer has
common factor, again up to a unit of \(\Q(t)\),
\[
 q(u)=10u_0^2+u_0u_1-4u_1^2.
\]
Its discriminant is \(161\), so it gives two distinct ruling lines
over the algebraic closure of \(\Q(t)\).  The exact generic-field
calculation also computes the saturation.  Writing
\(T=\left(\begin{smallmatrix}a&b\\c&d\end{smallmatrix}\right)\),
the saturated ruling ideal is
\[
\begin{split}
 I_{\mathrm{sat}}=(\delta,\,
 &10a^2+ac-4c^2,\,
 10ab+ad-4cd,\\
 &10b^2+bd-4d^2).
\end{split}
\]
The Groebner bases over \(\Q(t)\) give
\[
 H_{M_3}=(1,4,9,8,10,12,14,\ldots),\qquad
 H_{P/I_{\mathrm{sat}}}=(1,4,6,8,10,12,14,\ldots).
\]
The quotient \(I_{\mathrm{sat}}/I_{M_3}\) is therefore concentrated
in degree two and has length three.  This is the claimed embedded
vertex module, so the generic signature is \(R_2+E_3\).

For the two open-orbit rows \((1,2)\) and \((0,4)\), the previously
displayed base-point factors remain generic by orbit transport:
\[
 u_0u_1,\qquad 6u_0^2+u_0u_1+4u_1^2,
\]
and the second quadratic has nonzero discriminant.  Thus every
positive-dimensional rank-drop support in the table is certified on
its generic coefficient field, rather than inferred from an isolated
integral point.
The exact stabilizer ranks, slice coordinates, Groebner leading
monomials, denominator checks, and Hilbert counts are reproduced by
\texttt{checks/generic\_boundary\_atlas.py}.  Because each slice
map is dominant at its base point, the calculation is at the generic
point of the rank type rather than at a special witness.
\end{proof}

\subsection{Relative assassins and base change for the colon tower}

We prove the general bounded principal-colon theorem before applying it
to the determinant \(d\).  The order of the argument is essential.

\begin{lemma}[Finite geometric-assassin stratification]
\label{lem:finite-geometric-assassin-stratification}
Let \(X\to S\) be of finite presentation with \(S\) noetherian, and
let \(\mathcal F\) be a finitely presented coherent module which is
flat over \(S\).  After a finite locally closed stratification of
\(S\), there are finitely many closed support packets
\(Z_\nu\subset X\) such that for every geometric point
\(\bar s\to S\),
\[
 \Ass(\mathcal F_{\bar s})
 =
 \bigcup_\nu
 \{\text{generic points of irreducible components of }(Z_\nu)_{\bar s}\}.
\]
The packets may split geometrically.  After further finite refinement,
their dimensions, incidence relations, minimal-versus-embedded status,
and the numbers
\[
 \operatorname{length} H^0_{p\OO_{X_{\bar s},p}}(\mathcal F_{\bar s,p})
\]
at their associated generic points are constant.
\end{lemma}

\begin{proof}
Work componentwise on the base and use Noetherian induction.  Over a
geometric generic point the associated set is finite.  After a finite
extension of the generic field, collect Galois-conjugate associated
points into finitely many packets, take their closures in \(X\), and
shrink so that these closures are flat and their geometric fibre
dimensions are constant.  The base-change description of the relative
assassin in \cite[Tag 05AS]{StacksRelativeAssassin} permits this
geometric formulation.  For a flat module, the generic points of
fibres of the closure of a relative associated point remain relative
associated points after shrinking; we use
\cite[Tag 0GSJ]{StacksFlatRelativeAssassin}.

What remains is to exclude untracked fibrewise associated points.
The needed locus is constructible, not merely specialization closed.
We use the relative-assassin constructibility theorem
\cite[Tag 05KR]{StacksConstructibleAssassin}: for a quasi-compact open
\(U\subset X\), the locus of \(s\in S\) for which
\(\Ass(\mathcal F_s)\subset U_s\) is locally constructible.
Applying this to a finite affine cover adapted to the closures just
constructed, and intersecting the resulting constructible conditions,
gives a dense constructible locus on which the generic packets exhaust
the geometric fibrewise assassin.  A dense constructible subset of a
noetherian irreducible scheme contains a nonempty open.  Remove that
open and repeat on the noetherian complement.  This terminates and
gives a finite stratification.  The same constructibility argument
applied to pairwise closures fixes incidence and therefore
minimal-versus-embedded status.

For the multiplicity assertion, fix a tracked generic point \(p\).
The module
\[
 T_p=H^0_{p\OO_{X,p}}(\mathcal F_p)
\]
has finite length.  Choose \(n\) with \(p^nT_p=0\), spread both the
packet ideal and the finite torsion module after shrinking, and filter
it by the powers of that ideal.  The successive quotients are finite
modules on the packet.  Generic freeness makes their generic ranks
constant after shrinking, and the sum of those ranks is exactly the
displayed length.  Noetherian induction handles the complement.
\end{proof}

\begin{proof}[Proof of Theorem~\ref{thm:bounded-principal-colon-stratification}]
Fix \(0\le q<N\) and put
\[
 B=A/J,\qquad C_q=B/f^qB=A/(J,f^q),\qquad
 I_q=f^qB.
\]
There are two short exact sequences
\begin{equation}\label{eq:two-flat-sequences}
 0\longrightarrow I_q\longrightarrow B\longrightarrow C_q
 \longrightarrow0,
 \qquad
 0\longrightarrow K_q\longrightarrow A
 \xrightarrow{\ a\mapsto af^q\ } I_q\longrightarrow0.
\end{equation}
There is also
\[
 0\longrightarrow J\longrightarrow A\longrightarrow B\longrightarrow0.
\]
By generic flatness, simultaneously for the finite list
\(q=0,\ldots,N-1\), followed by Noetherian induction on the
complement, refine \(S\) so that \(B\) and every \(C_q\) are flat.
Then the first sequence in \eqref{eq:two-flat-sequences} shows that
every \(I_q\) is flat.  Since \(A\) is flat, the second sequence
remains exact after \emph{arbitrary} base change.  Since \(B\) is
flat, the sequence defining \(J\) does as well.

Let \(T\to S_\alpha\) be any base change.  Exactness identifies
\[
 B_T=A_T/J_T,\qquad
 (I_q)_T=f_T^qB_T,
\]
and the kernel of \(A_T\to(I_q)_T\) is consequently
\[
 (K_q)_T=(J_T:f_T^q).
\]
This is the desired colon base-change identity.  Notice that flatness
of the final cokernel \(C_q\) alone is not being used to preserve a
kernel; the image module \(I_q\) has first been proved flat.

Now refine once more so that every
\(M_q=A/(f,K_q)\), and the sources, targets and cokernels of the
finite list of multiplication maps under consideration, are flat.
Their Hilbert polynomials and fibre ranks are then constant and commute
with base change.  Apply
Lemma~\ref{lem:finite-geometric-assassin-stratification} to each
\(M_q\), and take a common finite refinement.  This produces the
support packets, incidences, minimal/embedded status and embedded
multiplicities asserted in the theorem.
\end{proof}

\begin{proof}[Proof of Proposition~\ref{prop:associated-prime-refinement}]
Apply Theorem~\ref{thm:bounded-principal-colon-stratification} to the
finite-type coefficient chart with
\[
 A=P_S,\qquad f=d,\qquad N=5.
\]
The hypothesis \(d^5\in J\) is
Theorem~\ref{thm:finite-all-corank-depth}, itself a specialization of
the universal symmetric-power determinant-completion theorem.
Include the finitely many graded multiplication maps in the same
refinement.  The theorem gives base change for every residual colon,
flatness and constant Hilbert polynomial for every \(\mathcal M_j\),
the exact geometric fibrewise associated-point packets with constant
embedded multiplicities, and constant multiplication ranks.  These
are precisely the three clauses of the proposition.

The construction of the coefficient strata is noncanonical, but its
fibrewise output is the intrinsic graded algebra, its geometric
assassin, and the canonical torsion lengths at associated generic
points.  Two presentation atlases therefore admit a common refinement
carrying the same primary signature.
\end{proof}

\subsection{Explicit orbit closures and graded specialization}

Let
\[
 e_1=h_1w_1,\quad e_2=h_1w_2,\quad
 e_3=h_2w_1,\quad e_4=h_2w_2
\]
be the standard basis of \(H\otimes W\).  The following families
are written before applying a general quadratic block \(\chi_t\).

\begin{theorem}[Corank-two orbit closure and specialization laws]
\label{thm:coranktwo-specialization-laws}
The closure relations needed to pass between all nine mixed-kernel
types are realized by the following one-parameter families:
\begin{enumerate}
\item rank three, off \(Q\) to on \(Q\):
\[
 \beta_t=
 \begin{pmatrix}0&1&0&0\\0&0&1&0\\-t&0&0&1\end{pmatrix},
 \qquad \ker\beta_t=\langle e_1+t e_4\rangle;
\]
\item secant to tangent:
\[
 \ker\beta_t=\langle e_1+e_4,\ e_2-t e_3\rangle,
\]
for which \(\delta|_{\Pj(\ker\beta_t)}=x^2+t y^2\);
\item secant to the two ruling families:
\[
 \langle e_1,e_2+t e_4\rangle\rightsquigarrow
 \langle e_1,e_2\rangle,\qquad
 \langle e_1,e_3+t e_4\rangle\rightsquigarrow
 \langle e_1,e_3\rangle;
\]
\item tangent to dual-rank-two:
\[
 \beta_t=
 \begin{pmatrix}0&1&-1&0\\0&0&0&t\\0&0&0&0\end{pmatrix};
\]
\item either ruling to dual-rank-one, by multiplying the second
nonzero row in the corresponding rank-two normal form by \(t\);
\item either rank-one type to the zero map, by replacing
\(\beta\) with \(t\beta\).
\end{enumerate}
For a general polynomial choice of \(\chi_t\), the generic fibre
and special fibre lie in the generic coefficient strata of the
indicated orbit types.  The corresponding primary signatures
specialize as follows:
\[
\begin{array}{l|l}
\text{degeneration}&\text{change of positive layers}\\ \hline
\text{off }Q\to\text{on }Q&V_1\to V_3\\
\text{secant}\to\text{tangent}&V_8\to V_8\\
\text{secant}\to H\text{-ruling}&V_8\to(V_{15},V_1)\\
\text{secant}\to W\text{-ruling}&V_8\to(R_1+E_5)\\
\text{tangent}\to\text{dual rank }2&V_8\to(R_2+E_3,V_1)\\
H\text{-ruling}\to\text{dual rank }1&(V_{15},V_1)\to(R_3,V_3)\\
W\text{-ruling}\to\text{dual rank }1&(R_1+E_5)\to(R_3,V_3)\\
\text{dual rank }2\to0&(R_2+E_3,V_1)\to(Q,Q)\\
\text{dual rank }1\to0&(R_3,V_3)\to(Q,Q).
\end{array}
\]
Thus ruling components and embedded vertex components are born in
actual one-parameter families of Fitting schemes.
\end{theorem}

\begin{proof}
The determinant of a tensor in each displayed kernel gives the stated
Segre incidence.  For example
\[
 \det\left(xI_2+y\begin{pmatrix}0&1\\-t&0\end{pmatrix}\right)
 =x^2+t y^2,
\]
which is secant for \(t\ne0\) and tangent for \(t=0\).  Likewise
\(\delta\) restricts to \(txy\) on each secant-to-ruling family.
The row-scaling families have the stated rank-one limits by inspection.

The residual maximal minors are polynomial in \(t\) and in the
coefficients of \(\chi_t\).  Each source and target generic stratum
is a nonempty open in its coefficient space.  Choosing
\(\chi_0\) in the target open and a general first-order perturbation
\(\chi_0+t\chi_1\) avoids the source bad divisors over
\(\C(t)\).  Hence the two fibres have the primary signatures in
Theorem~\ref{thm:generic-primary-signatures}.  Because the ideals
come from one polynomial family, the changes displayed in the table
are genuine specializations, not comparisons of unrelated normal
forms.
\end{proof}

There is an analogous pure-block degeneration.  In a fixed target
basis let
\[
 A_t=(f_1,f_2,t f_3):\Sym^2H\longrightarrow S_R
\]
and choose the remaining mixed and quadratic coefficients generally.
For \(t\ne0\), the induced mixed block is generically the rank-three
off-Segre type; at \(t=0\), \(A_H\) has rank two and the induced
mixed block has rank four.  Thus the first nontrivial layer specializes
from \(V_1\) to \(V_{22}\).  This gives an explicit
\(A_H\)-injective-to-rank-drop family, while the full table of
Theorem~\ref{thm:rankdrop-primary-signatures} describes the generic
pure-block boundary reached after further rank drops.
